\chapter{Podsumowanie i wnioski}\label{ch:podsumowanie}

\section{Podsumowanie zrealizowanych prac i celów projektu}\label{sec:podsumowanie-prac}

W ramach niniejszego projektu zaprojektowano, zaimplementowano oraz przetestowano kompletne rozwiązanie umożliwiające zautomatyzowane i bezpieczne zarządzanie regułami Zapory systemu Windows Defender w środowisku PowerShell 5.1. Opracowany system bezpośrednio rozwiązuje problemy zidentyfikowane podczas analizy standardowego interfejsu graficznego \file{wf.msc}:

\begin{enumerate}
    \item \textbf{Zastąpienie powolnego kreatora GUI zwalidowanym interfejsem CLI (Moduł~1)} - wieloetapowe okna dialogowe zastąpiono kreatorem konsolowym z natychmiastową kontrolą adresacji IPv4/IPv6, portów ($1\text{-}65535$) i protokołów oraz automatycznym kierunkowym mapowaniem portów (\code{LocalPort} dla \textit{Inbound}, \code{RemotePort} dla \textit{Outbound}). Konwencja \code{<NR>-OSO-<Action>-<Direction>-<HHmmss>} wyeliminowała problem przemieszania reguł własnych z setkami reguł systemowych.
    \item \textbf{Świadome zarządzanie kontekstem sieciowym (Moduł~2)} - zintegrowano podgląd reguł z diagnostyką usługi NLA, udostępniając informację o aktywnych kategoriach interfejsów sieciowych oraz elastyczny wybór profili (\textit{Domain}, \textit{Private}, \textit{Public}, \textit{Any}).
    \item \textbf{Transakcyjny audyt odporny na manipulacje (Moduł~3)} - zrealizowano silnik audytowy oparty na bezpośrednich migawkach rejestru i GPO w protokole zapisu wyprzedzającego (\code{intent} $\rightarrow$ \code{change} $\rightarrow$ \code{result}) z domyślną polityką \textit{fail-closed}, ochroną łańcuchem skrótów SHA-256 oraz zadaniem \textit{Watcher} wykrywającym i przypisującym zmiany zewnętrzne.
    \item \textbf{Ciągłość konfiguracji na zewnętrznym serwerze kopii (Moduł~4)} - wdrożono architekturę zdalnych kopii zapasowych z serwerem Ubuntu i usługą SFTPGo, pobierającą kopie \file{.wfw} w modelu \textit{pull} przed każdą zmianą reguł oraz weryfikującą skróty SHA-256 i kopię ratunkową \code{pre-restore} podczas odtwarzania.
\end{enumerate}

\section{Świadome ograniczenia rozwiązania i kierunki dalszego rozwoju}\label{sec:dalsze-prace}

Ocena każdego systemu bezpieczeństwa wymaga jawnego określenia granic przyjętego modelu zagrożeń. Ponieważ narzędzie działa na stacji roboczej kontrolowanej przez lokalnego administratora, zidentyfikowano następujące świadome ograniczenia architektoniczne oraz odpowiadające im mechanizmy kompensujące:

\begin{itemize}
    \item \textbf{Brak klucza tajnego w łańcuchu skrótów (SHA-256 bez HMAC)} - skrót każdego rekordu w \file{audit.jsonl} obliczany jest bezkluczową funkcją SHA-256, przez co lokalny administrator mógłby zmodyfikować historyczny wpis i przeliczyć kolejne skróty lub obciąć koniec pliku. Mechanizmem kompensującym jest eksport bieżącego nagłówka łańcucha (\textit{head hash}: \code{<seq>:<sha256>}, opcja menu \code{8} $\rightarrow$ \code{[H]}), którego zewnętrzne przechowanie i podanie podczas weryfikacji (opcja \code{6}) wykrywa zarówno obcięcie, jak i przepisanie dziennika.
    \item \textbf{Możliwość zatrzymania zadania \textit{Watcher} lub wyczyszczenia dziennika zdarzeń} - lokalny administrator może wyłączyć zadanie \code{\FirewallAudit\Watcher} lub wyczyścić kanał zdarzeń zapory. Działanie to jest jednak jawnie sygnalizowane: wyczyszczenie zdarzeń oznacza kolejne zmiany zewnętrzne jako pozbawione atrybucji (\code{unattributed}), a zatrzymanie zadania przełącza status w menu głównym w stan alarmowy \code{Stale} po 45~minutach.
    \item \textbf{Zależność od dostępności serwera kopii zapasowych} - domyślny przepływ wstrzymuje utworzenie lub usunięcie reguły w przypadku nieosiągalności serwera Ubuntu, chroniąc przed cichym rozbiegnięciem się stanu stacji z archiwum kopii. W sytuacjach awaryjnych administrator może świadomie pominąć wykonanie kopii zdalnej, zachowując pełną ciągłość pracy oraz lokalny zapis transakcyjny w dzienniku \file{audit.jsonl}.
\end{itemize}

Kierunki dalszego rozwoju systemu obejmują: (1) automatyczne przesyłanie każdego nowego nagłówka \textit{head hash} na serwer Ubuntu podczas wykonywania kopii zapasowej, (2) podpisanie skryptów certyfikatem Authenticode przy wymuszeniu polityki \code{AllSigned} oraz (3) rozszerzenie kreatora reguł o filtry ścieżek aplikacji (\code{-Program}) i nazw usług (\code{-Service}).

\section{Wnioski końcowe}\label{sec:wnioski}

Przeprowadzone prace projektowe i testy weryfikacyjne potwierdziły, że połączenie mechanizmów PowerShell oraz .NET z bezpośrednią analizą rejestru i zewnętrznym serwerem kopii zapasowych pozwala stworzyć narzędzie przewyższające standardowy interfejs graficzny szybkością działania i rygorem bezpieczeństwa. Architektura modułowa ze wspólnym kontraktem audytowym (\code{Invoke-AuditedChange}) oraz pobieraniem kopii w modelu \textit{pull} skutecznie chroni spójność i rozliczalność polityki zapory ogniowej.
